\section{A new quartic with 1367 variables}
\label{sec:new-quartic}

The finite-algebra construction minimizes the number of coordinates by
allowing a very high degree.  For a degree-four presentation, auxiliary
coordinates serve a different purpose: they keep each arithmetic relation
quadratic.  We now apply this method to the shared-point arithmetic system,
using the common height point, the fixed square scales, comparison order
$52$, scalar order $32$, and the semilocal size condition.  The scalar pole
in this quartic retains its preparation coordinate.  The separate
normalization which removes that coordinate in the smallest-variable
construction is not used here; its extra arithmetic would cost more
quadratic intermediate coordinates.

\subsection{The exact arithmetic conventions}

Represent a rational intermediate value by an affine linear form in the
coordinates already introduced.  Addition, subtraction, and multiplication
by a fixed rational coefficient then require no new coordinate.  A product
of two nonconstant affine forms $a,b$ is represented by one coordinate $v$
and the quadratic equation $v-ab=0$.  A rational inverse uses $av-1=0$.
More generally, an arbitrary quadratic expression $q$ can be assigned to
one new coordinate with $v-q=0$.  Common expressions are reused.

An element of $K=\mathbb Q(\sqrt2)$ is a pair of rational affine forms.
Its product with another pair is assigned to two coordinates, using
\[
 v_0-a_0b_0-2a_1b_1=0,\qquad
 v_1-a_0b_1-a_1b_0=0.
\]
This costs two coordinates, although an arithmetic algorithm minimizing
scalar multiplications might instead use three products and free additions.
The objective here is the number of existential coordinates.  An inverse
uses two coordinates and
\[
 a_0v_0+2a_1v_1=1,\qquad a_0v_1+a_1v_0=0.
\]
These equations hold exactly when the represented element is nonzero and
$(v_0,v_1)$ is its inverse.  They enforce the necessary domain without
introducing a separate norm and a separate inverse of that norm.

When the domain lemmas above already prove that $b\ne0$, a quotient
$z=a/b$ uses just the two equations
\[
 b_0z_0+2b_1z_1=a_0,\qquad b_0z_1+b_1z_0=a_1.
\]
Their determinant is $b_0^2-2b_1^2\ne0$, so their solution is unique.
This rule is used for doubling slopes, local parameters, logarithm ratios,
and the scalar local-parameter ratio, with the domain proof supplied in
the preceding subsection.  The subtraction slope retains an explicit
inverse of the difference of abscissas.  These domain arguments are part
of the proof of the quartic presentation; a coefficient-table check alone
would not justify dropping a denominator condition.

\subsection{Reusing polynomial evaluations}

Two elementary reorganizations save many auxiliary coordinates.  First,
write a polynomial with rational coefficients in blocks:
\[
 G(W)=\sum_{r=0}^s W^{br}
               \left(\sum_{j=0}^{b-1}c_{br+j}W^j\right).
\]
After computing $1,W,\ldots,W^b$, each parenthesized expression is an
affine linear combination of their two rational components.  Horner
evaluation in $W^b$ finishes the calculation.  This is the
Paterson--Stockmeyer principle \cite{PatersonStockmeyer1973}.  Here $b=3$
for $G_{12}$ and $b=4$ for $G_3$.  The identity is obtained simply by
grouping the coefficients according to their exponent modulo $b$; no
approximation of the logarithm is introduced by this evaluation rule.

Second, the eighteen field tests $D_j$ are the same at every coordinate
detector.  If $\delta=A\lambda-1$, the corresponding factor is
\[
 1+q_mD_j^2,\qquad q_m=R^{1-2m}\delta^{-2m-2},
 \qquad m\in\{1,16,52\}.
\]
For each order $m$, form once the polynomial
\[
 E_m(U)=\prod_{j:m_j=m}(1+UD_j^2)\in K[U].
\]
Its coefficients are the elementary symmetric polynomials in the listed
$D_j^2$.  There are respectively ten, two, and six factors.  For each of
the nine detectors, evaluate $E_1(q_1)E_{16}(q_{16})E_{52}(q_{52})$.
This is exactly the product defining $W$, with the field coefficients
reused across all nine detectors.

Finally, for $W=w_0+w_1\sqrt2$,
\[
 \operatorname{Tr}(W^{-1})=\frac{2w_0}{w_0^2-2w_1^2},\qquad
 B=\frac{R(w_0^2-2w_1^2)^2}{4w_0^2}.
\]
The inverse of $w_0$ enforces the trace condition.  Since $w_0\ne0$,
the norm of $W$ is nonzero in the quadratic field.  Together with
$R\ne0$, this also gives $B\ne0$, so an additional inverse of $B$
would be redundant.

\subsection{Finite choices at degree four}

The point occurrences use four rational halfpoint coordinates and two
curve equations.  The three-coordinate quartic chart is more economical
for finite-algebra elimination, but is not needed in this degree-four
construction.  The nine point occurrences are the four root points and
the five shared comparison points; the first comparison point is also
the scalar-height point.

The fixed square scales proved above remove all zero-component case
choices.  Only the scalar correction remains.  For each
$p\in\{2,3,5,7,11,13,17\}$ introduce a Boolean coordinate $\epsilon_p$,
and introduce one sign coordinate $\epsilon_-$.  Impose
$\epsilon^2-\epsilon=0$ for all eight coordinates and compute
\[
 c=(1-2\epsilon_-)\prod_p\bigl(1+(p-1)\epsilon_p\bigr).
\]
Over $\mathbb Q$, this realizes exactly the required 256 signed squarefree
corrections.  Unlike a finite product of alternative polynomials, this
keeps the total degree four after quadratic compilation, at the cost of
eight coordinates.

\subsection{Removing redundant quadratic coordinates}

The compiler next performs two reversible operations.  Replacing a
residual $q_i$ by $q_i+cq_j$, while retaining $q_j$, does not change the
common zero set.  If a residual has the form
\[
                  av+q=0,\qquad a\in\mathbb Q^\times,
                  \quad v\notin q,
\]
we may substitute $v=-q/a$ in all the other residuals and remove this
coordinate and equation.  Every zero of the new system reconstructs one
and only one value of $v$; the converse follows by substitution.  The
compiler accepts such a substitution only if every resulting residual
still has degree at most two.  In particular a quadratic replacement is
permitted only when that coordinate occurs linearly elsewhere.  Its exact
row operations and substitutions are retained as a replayable certificate.

After these operations, clear the rational denominator of each residual
and divide its integer coefficients by their common divisor.  Neither
operation changes its rational zero set.  If the resulting residuals are
$q_1,\ldots,q_s$, define
\[
                         P=\sum_{i=1}^s q_i^2.
\]
The rational zero set of $P$ is exactly their common zero set, because
each summand is nonnegative.  All these equivalences concern rational
solutions and do not require numerical sampling.

\begin{theorem}[The supplied quartics]\label{thm:quartic-1367}
Assume the source arithmetic inputs and Sun's simultaneous prime-pattern
theorem used in the shared-point construction.  The supplied polynomial
$P_{1367}$ has 1367 total variables, 1366 existential witnesses, and total
degree four, and defines $\mathbb Z$ in $\mathbb Q$.  It is specified by
a complete sparse coefficient table over $\mathbb Z[\kappa]$, where
\[
                         \kappa=3^{2^{28}}
\]
is a fixed integer, not a coordinate.  There is also a fully expanded
ordinary integer quartic $P_{1391}$ with 1391 total variables and primitive
coefficient height less than $2^{1020}$, with the same defining property.
\end{theorem}
\begin{proof}
Apply the just-described rules to the shared-point source with the
precision and size choices specified at the start of this section.
Each arithmetic coordinate has its prescribed value, except where the
explicit finite Boolean choices are intended; the paired quotient domain
lemmas establish the same assertion for the divisions.  The reversible
quadratic operations preserve the rational input projection.  The source
arithmetic theorem therefore identifies the final projection with
$\mathbb Z$.

The exact implementation introduces 1367 coordinates after sixty
quadratic-preserving eliminations.  Its 1322 primitive residuals expand
to 29,419 nonzero terms over $\mathbb Z[\kappa]$.  The independent checker
replays every arithmetic relation, row operation, and elimination and
independently convolves every pair of residual terms.  This checks the
entire sparse table, including mixed coefficients.  Its maximum variable
degree is exactly four.

For the ordinary integer alternative, start with the fixed integer
$3^{16}$ and use twenty-four successive squarings to compute
$3^{16\cdot2^{24}}=3^{2^{28}}$.  The associated coordinates are retained
to avoid expanding this integer into the coefficients.  Their defining
equations have degree at most two and uniquely determine their values.
The resulting 1391-coordinate system has 1346 residuals and 29,489
nonzero expanded integer terms.  The independent expansion check gives
coefficient bit length at most 1020.  The coefficient content is one.
These facts prove the degree, height, and count assertions.
\end{proof}

The supplied data are in
\nolinkurl{research/phase2/quartic/quartic-best/polynomial.json.gz} and
\nolinkurl{research/phase2/quartic/quartic-small-coefficients/polynomial.json.gz}.
An entry $[\mathbf m,k,c]$ means $c\kappa^k\prod_{i\in\mathbf m}X_i$,
with repetitions in $\mathbf m$ recording powers.  The index zero is the
input.  In the ordinary integer table every $k$ is zero.  This convention,
the complete tables, and the exact construction give ordinary finite
integer polynomials.  No gate or coefficient symbol is silently omitted
from the variable count.

The program \nolinkurl{build_shared_quartic.py} is a producer; it is not
trusted by the verification program
\nolinkurl{verify_shared_quartic.py}, which does not import it.
The separate mathematical review checks the source formulas and the
domain lemmas.  Neither computational check is claimed to verify the
external arithmetic theorems.  No global minimality claim is made for
1367 variables: the theorem supplies a concrete upper bound for the
quartic problem, while the unconditional lower bound for this
sum-of-squares architecture is five total variables.
